\cleardoublepage
\phantomsection
\pdfbookmark{Summary}{Summary}

\chapter*{Abstract}
\markboth{Summary}{Summary}

\begingroup
\setlength{\parindent}{0pt}
\setlength{\parskip}{0.8em}

Multimodal artificial intelligence has advanced significantly in recent years, particularly in mapping visual and textual data into shared latent spaces. However, extending this success to three-dimensional geometric representations has proven exceptionally challenging. 3D point clouds possess topologically dissimilar and non-Euclidean structures that fundamentally differ from sequential language models; this geometric mismatch renders zero-shot 3D-text alignment a critical open problem.

This Master's thesis addresses this challenge directly through a comprehensive empirical and theoretical investigation of cross-modal latent alignment between uni-modal 3D objects and text. We extract features from a large-scale dataset of 3D objects and their corresponding textual descriptions utilizing several pre-trained neural architectures --- PointNet, OpenCLIP, \gls{roberta}, and \gls{bert} --- treating each modality as a high-dimensional global embedding.

We begin by rigorously replicating and verifying the state-of-the-art supervised linear alignment framework proposed by \textcite{hadgi2025escaping}. Anchor pairs are projected into a shared, low-dimensional space via an exact \gls{svd}-based \gls{cca}, followed by an affine least-squares mapping, establishing a robust performance benchmark. Evaluated on a held-out query set and averaged over multiple random seeds, our replication confirms that supervised linear transformations effectively expose shared geometric structure between the two modalities --- achieving a $35.7\%$ Top-5 retrieval accuracy on the PointNet $\times$ \gls{clip} pairing.

Subsequently, we evaluate unsupervised and weakly-supervised \gls{ot}-based alignment strategies, leveraging \gls{gw} distance metrics \cite{vayer2019optimal}. We test five GW variants --- \gls{sgw}, \gls{risgw}, \gls{ssrisgw}, \gls{fgw}, and \gls{lrgw} --- integrated with the intra-domain scale-matching normalisation introduced by \textcite{li2026scale}. These unsupervised formulations uniformly fail to achieve meaningful alignment. Across thirteen independent ablation experiments, every unsupervised GW approach collapses, yielding retrieval rates statistically indistinguishable from random chance across all tracked metrics. As an original contribution, we also implement and evaluate two supervised non-linear extensions: \gls{kcca} with Nystr\"{o}m approximation, which reaches $15.8\%$ Top-5, and \gls{rml}, which performs considerably better at $27.2\%$ --- capturing approximately 76\% of the supervised linear baseline's performance.

Finally, we establish the theoretical basis for this alignment failure through a detailed structural analysis. By analyzing transport plan entropy, mass concentration ratios, and self-structural GW cost ratios, we demonstrate that 3D point cloud spaces and text latent spaces fundamentally violate the assumption of structural isometry. Because the intra-domain distance distributions are not isometrically correlated, the quadratic solvers driving unsupervised optimal transport are mathematically forced to converge toward non-informative, high-entropy local minima. Ultimately, this proves that explicit anchor supervision is not merely beneficial, but strictly necessary for successful cross-modal 3D-text alignment --- establishing a definitive mathematical boundary on the capabilities of relational optimal transport in multimodal learning.

\endgroup
\vfill